% !TeX root = propuesta_prueba_semana04_curso_2.tex
% !TeX program = pdflatex
\documentclass[10pt,a5paper]{article}
\newif\ifmostrarSoluciones
\mostrarSolucionestrue
\usepackage{estilo_apuntes}
\encabezado{Semana 4}{Primera evaluación -- versión B}

\begin{document}
\tituloclase{Primera evaluación}{Propuesta para el segundo curso}{Semanas 1 a 3}

\noindent
\textbf{Nombre:} \rule{7cm}{0.4pt}

\medskip
\noindent
\textbf{Fecha:} \rule{3cm}{0.4pt}
\hfill
\textbf{Puntaje:} \rule{1.5cm}{0.4pt}

\begin{nota}[indicaciones]
\begin{itemize}
  \item La duración de la prueba es de 1 hora y 30 minutos.
  \item Escriba el desarrollo de cada ejercicio de manera clara y ordenada.
  Justifique cada paso e indique las propiedades, restricciones o resultados que
  utilice. Una respuesta sin justificación recibirá únicamente crédito parcial.
  \item No se permite el uso de ningún dispositivo electrónico durante la
  evaluación.
  \item Cualquier intento de copia penalizará la evaluación con una calificación
  total de cero.
  \item Puntaje sugerido: 10 puntos por ejercicio.
\end{itemize}
\end{nota}

\begin{ejercicios}[propuesta paralela]
\begin{enumerate}
  \item Para $a\in\mathbb{R}$, considere la función
  \[
    F_a\colon\mathbb{R}\to\mathbb{R},
    \qquad
    F_a(x)=
    \begin{cases}
      \dfrac{x^2+ax-2a-4}{|x-2|}, & x\neq2,\\[2mm]
      7, & x=2.
    \end{cases}
  \]
  Calcule los límites laterales de $F_a$ cuando $x$ tiende a $2$ y determine
  todos los valores de $a$ para los cuales existe
  \[
    \lim_{x\to2}F_a(x).
  \]
  Para esos valores, calcule el límite y compárelo con $F_a(2)$.

  \item Sea $A$ una función definida cerca de $0$, excepto posiblemente en
  $0$, que satisface
  \[
    |A(x)|\leq3.
  \]
  Sin suponer que $A$ tiene límite, calcule mediante el teorema del Sanduche
  \[
    \lim_{x\to0}
    \frac{x^2}{\sqrt{x^2}+x^2}\,A(x).
  \]

  \item Considere las funciones
  \[
    f\colon\mathbb{R}\setminus\{-1\}
    \to\mathbb{R}\setminus\{2\},
    \qquad
    f(x)=\frac{2x+1}{x+1},
  \]
  y
  \[
    g\colon\mathbb{R}\setminus\{2\}
    \to\mathbb{R}\setminus\{-1\},
    \qquad
    g(x)=\frac{1-x}{x-2}.
  \]
  \begin{enumerate}
    \item Verifique que las leyes de asignación producen valores en los
    codominios indicados.
    \item Determine si cada función es inyectiva, sobreyectiva o biyectiva.
    \item Halle la ley de asignación y el dominio de $g\circ f$ y de
    $f\circ g$. Decida si $f$ y $g$ son funciones inversas.
  \end{enumerate}
\end{enumerate}
\end{ejercicios}

\ifmostrarSoluciones
\begin{solucion}[de la propuesta paralela]
\begin{enumerate}
  \item El numerador se factoriza como
  \[
    x^2+ax-2a-4=(x-2)(x+a+2).
  \]
  Si $x<2$, entonces $|x-2|=-(x-2)$ y
  \[
    F_a(x)=-(x+a+2).
  \]
  Si $x>2$, entonces $|x-2|=x-2$ y
  \[
    F_a(x)=x+a+2.
  \]
  Por tanto,
  \[
    \lim_{x\to2^-}F_a(x)=-(a+4),
    \qquad
    \lim_{x\to2^+}F_a(x)=a+4.
  \]
  El límite bilateral existe si y solo si
  \[
    -(a+4)=a+4,
  \]
  de donde $a=-4$. En ese caso,
  \[
    \lim_{x\to2}F_{-4}(x)=0,
    \qquad
    F_{-4}(2)=7.
  \]
  El límite y el valor de la función en el punto son diferentes.

  \item Para $x\neq0$, se tiene $\sqrt{x^2}=|x|$. Entonces
  \[
    \left|
      \frac{x^2}{\sqrt{x^2}+x^2}A(x)
    \right|
    \leq
    3\frac{x^2}{|x|+x^2}
    =3\frac{|x|}{1+|x|}
    \leq3|x|.
  \]
  En consecuencia,
  \[
    -3|x|
    \leq
    \frac{x^2}{\sqrt{x^2}+x^2}A(x)
    \leq
    3|x|.
  \]
  Como ambas cotas tienden a $0$, el teorema del Sanduche permite concluir que
  \[
    \lim_{x\to0}
    \frac{x^2}{\sqrt{x^2}+x^2}A(x)=0.
  \]

  \item Si $f(x)=2$, entonces $2x+1=2(x+1)$, lo cual implicaría $1=2$.
  Por tanto, $f(x)\neq2$. De manera análoga, si $g(x)=-1$, se obtendría
  $1-x=-(x-2)$, que también implicaría $1=2$. Las dos leyes producen valores
  en sus codominios.

  Para estudiar $f$, sea $y\in\mathbb{R}\setminus\{2\}$. Al resolver
  \[
    y=\frac{2x+1}{x+1}
  \]
  para $x$, se obtiene el único valor
  \[
    x=\frac{1-y}{y-2}.
  \]
  Este valor no es $-1$, por lo que pertenece al dominio de $f$. Así, cada
  elemento del codominio tiene una única preimagen y $f$ es biyectiva. El mismo
  argumento, o la simetría de las expresiones, muestra que $g$ es biyectiva.

  Para todo $x\in\mathbb{R}\setminus\{-1\}$,
  \begin{align*}
    (g\circ f)(x)
    &=\frac{1-\dfrac{2x+1}{x+1}}
            {\dfrac{2x+1}{x+1}-2}\\
    &=x.
  \end{align*}
  Por tanto,
  \[
    \operatorname{Dom}(g\circ f)=\mathbb{R}\setminus\{-1\}.
  \]
  Para todo $x\in\mathbb{R}\setminus\{2\}$,
  \begin{align*}
    (f\circ g)(x)
    &=\frac{2\dfrac{1-x}{x-2}+1}
            {\dfrac{1-x}{x-2}+1}\\
    &=x.
  \end{align*}
  Por tanto,
  \[
    \operatorname{Dom}(f\circ g)=\mathbb{R}\setminus\{2\}.
  \]
  Las dos composiciones son las identidades en los dominios correspondientes;
  en consecuencia, $f$ y $g$ son funciones inversas.
\end{enumerate}
\end{solucion}
\fi

\end{document}
